% 使用 XeLaTeX 编译。匿名初稿保留封面身份占位符。
\documentclass{gmcmthesis}
\usepackage[section]{placeins}
\usepackage{needspace}
\usepackage{subfig}
\usepackage{float}
\usepackage{longtable}
\usepackage{array}
\usepackage{siunitx}
\setcitestyle{super,square,comma}
\title{山区洪涝灾害下无人机运输与通信协同优化}
\baominghao{[参赛队号]}
\schoolname{[学校名称]}
\membera{[成员A]}
\memberb{[成员B]}
\memberc{[成员C]}
\begin{document}
\maketitle
\begin{abstract}
针对山区洪涝灾害中道路受阻、应急物资需求分散且通信链路受地形遮挡的问题，本文围绕无人机运输、时限排程、空中中继和任务分区建立分层优化框架。问题一依据无人机载荷与能耗约束确定服务区安全载荷，并将不可拆货箱划分为运输批次；参考结果给出 15 个服务区、80 个货箱的单区组批方案，共 18 架次，运输能耗为 57.937\,kWh。问题二在批次基础上安排运输机与电池资源，得到 20 个架次，其中 2 个架次串联服务两个服务区；参考排程完工时刻为 9980\,s，总能耗为 61.208\,kWh。问题三按参考解的通信模式配置直连与中继，20 个运输架次中 14 个采用直连、6 个采用中继，对应 6 项中继任务，报告中继能耗为 1.797\,kWh。问题四将运输架次连接的服务区视为不可拆组件，对 K=2、3 的独立任务分区进行峰值资源核算；根据用户提供的更新版 Q4 TEX，两个方案的最低库存缺口分别为 3 和 4 个单位，K=2 较少。

本文所有数值均取自工作区归档参考结果文件或用户提供的新版 Q4 TEX；本稿未重新求解，也未对参考解作独立审计，因而结果应视为初稿报告值。文中着重说明分层建模关系、资源约束及结论适用范围。
\keywords{应急物流\quad 无人机调度\quad 空地协同通信\quad 资源分区}
\end{abstract}
\pagestyle{plain}
\section{问题重述}
\subsection{问题背景}
洪涝灾害可能造成道路中断，使常规运输难以在首批救援窗口内到达偏远村落。无人机可绕开地面通行障碍，但其载荷、续航、电池周转和通信距离均有限；山区地形起伏还会引起视距遮挡。因此，方案设计需要同时考虑物资批次、运输时序、通信保障及多组独立执行时的资源配置。

\subsection{问题内容}
题目给定应急物资货箱、服务区、调度中心、运输无人机及电池参数，并给出地形和通信数据。本文按题意将任务拆为四个相互衔接的问题：
\begin{enumerate}
  \item 对各服务区计算不同机型的安全载荷，并将不可拆货箱组成满足载荷、容积和能源约束的单区运输批次。
  \item 在运输批次基础上确定多区路线、起飞时间、无人机及电池分配，使需求获得服务并满足硬时限与资源周转条件。
  \item 对运输过程中的地空链路进行通信模式配置；在直连不足时引入空中中继，并统计中继任务与额外能源需求。
  \item 冻结既有运输和通信任务，将服务区分为 $K=2$ 或 $K=3$ 个独立任务组，保证同一跨区运输任务不被拆分，比较各组峰值并发资源与库存缺口。
\end{enumerate}

\subsection{技术路线}
本文采用“安全载荷与组批—时限排程—通信模式选择—固定方案分区核算”的递进式路线。前一问题输出作为后一问题的固定输入，使每问新增决策边界清晰。全流程示意见图~\ref{fig:route}。
\begin{figure}[htbp]
  \centering\includegraphics[width=0.96\linewidth]{technical_route.png}
  \caption{四问协同技术路线}\label{fig:route}
\end{figure}

\section{问题分析}
\subsection{问题一：单区组批}
本问的关键是将连续载荷能力转化为离散货箱批次。各机型对每个服务区的往返能力不同，且货箱不可拆分，故需先获得机型—服务区安全载荷，再在质量和容积限制下进行组合。结果要同时报告批次数和能耗，避免只按质量容量而忽略续航。

\subsection{问题二：时限排程}
组批确定后，问题转化为带时间窗的异构资源排程。每项运输任务占用一架无人机和一组对应电池；飞行结束后无人机返场，电池还需充电才能再次使用。双服务区路线可减少重复起降，但会延长任务时间并改变逐箱送达时刻，因此需将路线、起飞时刻和资源占用统一表示。

\subsection{问题三：通信保障}
通信配置继承运输轨迹。地面站—运输机直连可用时不额外安排中继；若直连无法支撑全程通信，则用运输机—中继机与中继机—地面站两跳链路覆盖缺口。中继任务包含转场与服务时间，并消耗中继机和能源组件。

\subsection{问题四：固定任务分区}
跨服务区运输架次形成不可拆分关系，可先构造服务区连接图并取连通分量，再将分量分配给独立任务组。分区之后，各组不能共享库存，资源总需求由组内并发占用的峰值决定；因此比较对象应为峰值并发需求相对库存的缺口，而不是简单累加各区任务数。

\section{模型假设}
在题目给定参数与参考排程口径下，采用以下简化假设：
\begin{enumerate}
  \item 货箱作为不可拆分任务处理，其质量、体积、类别和时限按输入记录取值；本文不对输入数据作统计插补。
  \item 同一运输架次按给定闭环路线完成服务，架次开始、返航及逐箱送达时刻采用参考结果记录的时间口径。
  \item 无人机在一次运输任务中连续占用，电池从任务开始占用至达到下一次可用时刻；充电周转由结果文件中的可用时刻表示。
  \item 通信模式和中继任务沿用参考结果，不在本文中重新搜索候选位置或重优化通信链路。
  \item 问题四冻结运输和通信任务时序；各分区独立执行且不跨组共享无人机、电池或中继资源。
  \item 峰值资源需求按任务区间并发数计量。该指标表示固定时序下的并发配置下界，不等同于实体设备 ID 使用数量，也不单独刻画组间转场。
\end{enumerate}
上述假设用于明确参考结果的适用范围；风雨扰动、临时禁飞区、链路随机衰落和资源故障不在本稿数值范围内。

\section{符号说明}
\begin{table}[htbp]
\centering
\caption{主要符号说明}
\begin{tabular}{c p{8.5cm} c}
\toprule
符号 & 含义 & 单位\\
\midrule
$i\in\mathcal I$ & 服务区索引，$|\mathcal I|=15$ & --\\
$b\in\mathcal B$ & 不可拆分货箱索引，$|\mathcal B|=80$ & --\\
$k\in\mathcal K$ & 无人机机型，$\mathcal K=\{A,B,C\}$ & --\\
$r\in\mathcal R$ & 可分配资源类型（运输机、电池、中继机或中继组件） & --\\
$P_{ik}$ & 机型 $k$ 在服务区 $i$ 的安全载荷 & kg\\
$x_{bj}$ & 货箱 $b$ 是否分配至运输批次 $j$ & 0--1\\
$s_j,e_j$ & 架次 $j$ 的开始与结束时刻 & s\\
$E_j$ & 架次 $j$ 的运输能耗 & kWh\\
$g_{ih}$ & 服务区组件 $i$ 是否分配至任务组 $h$ & 0--1\\
$N_{hr}^{\min}$ & 任务组 $h$ 对资源 $r$ 的峰值并发需求 & 个\\
$S_r$ & 资源 $r$ 的现有库存 & 个\\
$\Delta_r$ & 资源 $r$ 的库存缺口 & 个\\
\bottomrule
\end{tabular}
\end{table}
符号随各问题的局部下标含义在首次使用处补充说明。

\section{问题一：安全载荷与单区组批}
\subsection{安全载荷约束}
对服务区 $i$ 与机型 $k$，记往返飞行能量为 $E_{ik}(m)$，电池可用能量为 $B_k$，最低返航储备比例为 $\rho$。安全载荷 $P_{ik}$ 由下式定义：
\begin{equation}
P_{ik}=\max\{m: E_{ik}(m)\leq (1-\rho)B_k,\;m\leq \bar P_k\},\qquad
0\leq P_{ik}\leq \bar P_k.
\label{eq:q1payload}
\end{equation}
其中 $\bar P_k$ 为机型额定载荷上限。参考结果中各区安全载荷主要受额定上限约束，少数远距离/地形条件下的 C 型上限低于 80\,kg；S008 的 B、C 型安全载荷分别为 29.501\,kg 和 59.945\,kg。

\subsection{不可拆货箱组批}
令 $x_{bj}=1$ 表示货箱 $b$ 分配给架次 $j$，$y_{jk}=1$ 表示架次 $j$ 选用机型 $k$。对每箱唯一分配、区内批次一致服务区、质量/容积和机型选择约束写为
\begin{align}
&\sum_j x_{bj}=1 &&(\forall b),\\
&\sum_{b\in\mathcal B_i}m_bx_{bj}\leq \sum_k P_{ik}y_{jk} &&(\forall i,j),\\
&\sum_{b\in\mathcal B_i}v_bx_{bj}\leq \sum_k V_ky_{jk} &&(\forall i,j),\\
&\sum_k y_{jk}\leq 1 &&(\forall j),\qquad x_{bj},y_{jk}\in\{0,1\}.
\end{align}
在可行分配中以架次规模及运输能耗作为参考解的比较指标。本文不在本稿重新运行该组批算法，以下数值均转录自归档参考结果。

\subsection{组批结果}
15 个服务区共形成 18 个单区往返架次，其中 B 型和 C 型各 9 架次；累计运输能耗为 57.937\,kWh。图~\ref{fig:q1payload} 汇总机型—服务区安全载荷，图~\ref{fig:q1type} 与图~\ref{fig:q1energy} 分别给出机型架次构成和服务区能耗。可见单一机型并不能覆盖所有载荷与航程组合，安全载荷应随服务区分别计算。
\begin{figure}[htbp]\centering\includegraphics[width=0.82\linewidth]{fig_q1_payload_heatmap.png}\caption{各区机型安全载荷}\label{fig:q1payload}\end{figure}
\begin{figure}[htbp]\centering\includegraphics[width=0.62\linewidth]{fig_q1_type_count.png}\caption{问题一机型架次}\label{fig:q1type}\end{figure}
\begin{figure}[htbp]\centering\includegraphics[width=0.92\linewidth]{fig_q1_area_energy.png}\caption{问题一分区能耗}\label{fig:q1energy}\end{figure}

\subsection{返航储备敏感性}
改变返航储备比例 $\rho$，比较 15 区平均安全载荷。图~\ref{fig:q1sens} 显示储备比例由 10\% 增至 30\% 时，C 型平均安全载荷由约 79.08\,kg 降至 69.59\,kg，B 型由 30.00\,kg 降至 28.68\,kg，A 型平均值由 25.00\,kg 降至 24.40\,kg。C 型变化幅度较大，说明重载机型的可用载荷对返航储备设定更敏感。
\begin{figure}[htbp]\centering\includegraphics[width=0.88\linewidth]{fig_q1_reserve_sensitivity.png}\caption{返航储备敏感性}\label{fig:q1sens}\end{figure}

\section{问题二：多区运输与资源排程}
\subsection{时序决策模型}
令 $z_{jr}=1$ 表示运输任务 $j$ 使用资源 $r$，$s_j$ 为起飞时刻，$e_j=s_j+d_j$ 为返航时刻。货箱 $b$ 在任务 $j$ 中的送达偏移为 $\delta_{bj}$，则送达时刻 $t_b=s_j+\delta_{bj}$。每箱唯一服务、货箱硬时限及架次与机型适配关系为
\begin{align}
&\sum_{j:b\in j}1=1,\qquad t_b\leq D_b &&(b\in\mathcal B_{\rm hard}),\\
&\sum_{r\in\mathcal R(j)}z_{jr}=1 &&(\forall j),\\
&[s_j,e_j)\cap[s_\ell,e_\ell)=\emptyset &&\text{若任务 }j,\ell\text{ 使用同一架无人机},\\
&[s_j,c_j)\cap[s_\ell,c_\ell)=\emptyset &&\text{若任务 }j,\ell\text{ 使用同一电池}.
\end{align}
其中 $D_b$ 是硬时限，$c_j$ 是电池重新可用时刻。对多区路径，送达偏移随访问次序计算，且同一架次货箱共享路线。参考排程报告为 20 个运输架次、完工时刻 9980\,s、运输能耗 61.208\,kWh。

\subsection{排程结构与资源占用}
图~\ref{fig:q2gantt} 按实体运输无人机展示参考架次占用区间。图~\ref{fig:q2energy} 按机型汇总能耗，图~\ref{fig:q2ecdf} 给出逐箱交付时刻的经验累计分布，图~\ref{fig:q2mix} 对比机型架次和能耗。参考解中有 2 个架次访问两个服务区，其余为单区任务。图示用于呈现参考排程结构，不代表本稿重新计算或独立验证。
\begin{figure}[htbp]\centering\includegraphics[width=0.95\linewidth]{fig_q2_gantt.png}\caption{运输无人机任务时序}\label{fig:q2gantt}\end{figure}
\begin{figure}[htbp]\centering\includegraphics[width=0.72\linewidth]{fig_q2_type_energy.png}\caption{问题二机型能耗}\label{fig:q2energy}\end{figure}
\begin{figure}[htbp]\centering\includegraphics[width=0.84\linewidth]{fig_q2_delivery_ecdf.png}\caption{逐箱交付累计曲线}\label{fig:q2ecdf}\end{figure}
\begin{figure}[htbp]\centering\includegraphics[width=0.78\linewidth]{fig_q2_aircraft_mix.png}\caption{问题二机型结构}\label{fig:q2mix}\end{figure}
参考结果随附的 verification JSON 报告了货箱覆盖、载荷、硬时限、设备与电池区间和能耗检查状态；这里仅记录其报告值，不将该文件视为本轮独立审计。实际应用中还应结合天气、临时禁飞区和充电功率波动滚动更新排程。

\section{问题三：运输任务的通信保障}
\subsection{通信模式与中继任务}
令 $u_j=1$ 表示运输任务 $j$ 采用中继，$u_j=0$ 表示直连。以候选通信方案 $p\in\mathcal P_j$ 为单位，$w_{jp}$ 表示任务选择的模式、位置及资源组合。每个任务选择一个候选方案，并满足运输机通信模式连续覆盖的参考关系：
\begin{equation}
\sum_{p\in\mathcal P_j}w_{jp}=1,\qquad w_{jp}\in\{0,1\}.
\end{equation}
候选链路的可用性依据地形视距和双向链路预算确定；中继模式还需覆盖中继机转场、服务和返航区间，并分配相应的能源组件。该节沿用归档参考解中的候选选项与任务，不重新生成链路判定或优化通信位置。

\subsection{参考结果}
归档结果列出 64 个候选选项关系，最终运输模式计数为 14 个直连、6 个中继；配置 6 项中继任务，中继能耗合计 1.797\,kWh。图~\ref{fig:q3mode} 表示两种模式的架次计数，图~\ref{fig:q3relayenergy} 展示各中继任务的能耗。中继任务能耗约为 0.218--0.367\,kWh，差异来自任务转场与服务时长的不同。
\begin{figure}[htbp]\centering\includegraphics[width=0.66\linewidth]{fig_q3_mode_count.png}\caption{通信模式架次构成}\label{fig:q3mode}\end{figure}
\begin{figure}[htbp]\centering\includegraphics[width=0.86\linewidth]{fig_q3_relay_energy.png}\caption{中继任务能耗}\label{fig:q3relayenergy}\end{figure}

\subsection{适用范围}
以上数据用于描述给定参考解所采用的通信模式、任务数和能源统计口径。本文不将其表述为对所有通信资源并发冲突的独立审计结论，也不据此宣称联合全局最优。若部署要求不同中继机任务必须由单架实体机串行完成，应另行采用实体机区间约束并更新排程和结果。

\section{问题四：固定任务分区与资源缺口}
\subsection{不可拆组件与分区变量}
问题四保持问题二的运输架次、问题三的通信任务及其时序不变，只对服务区进行独立任务分区。若同一运输架次涉及多个服务区，这些服务区必须归入同一组。以运输架次构造服务区无向连接图，其连通分量为不可拆组件。对 15 个服务区得到 13 个组件，其中 $\{S008,S011\}$ 与 $\{S010,S013\}$ 分别由双服务区架次连接，其余组件各含一个服务区。

令 $g_{ih}\in\{0,1\}$ 表示组件 $i$ 是否分配到任务组 $h$，$h=1,\ldots,K$，其中 $K\in\{2,3\}$。每个组件恰分配一次且每组非空：
\begin{equation}
  \sum_{h=1}^{K}g_{ih}=1\quad(\forall i),\qquad
  \sum_i g_{ih}\geq 1\quad(\forall h).
  \label{eq:q4assign}
\end{equation}
以组件为分配对象可直接保证同一架次的服务区不可拆分。

\subsection{组内峰值资源需求}
固定任务开始和结束时刻。记 $I_{jr}$ 为任务 $j$ 对资源类型 $r$ 的占用区间，$\mathbf{1}(\cdot)$ 为示性函数，则组 $h$ 对资源 $r$ 的最少并发需求为
\begin{equation}
  N_{hr}^{\min}=\max_t\sum_{j\in h}\mathbf{1}(t\in I_{jr}).
  \label{eq:q4peak}
\end{equation}
运输无人机区间取起飞至返航；电池区间延伸至充电结束；中继无人机区间包括出发、建链、服务和返航；中继能源组件区间还覆盖充电周转。对库存量 $S_r$，分区后的配置缺口为
\begin{equation}
  \Delta_r=\max\left(0,\sum_{h=1}^{K}N_{hr}^{\min}-S_r\right),\qquad
  \Delta=\sum_r\Delta_r.
  \label{eq:q4gap}
\end{equation}
根据用户提供的更新版 Q4 TEX，对 13 个组件的全部非空 $K$ 分区规范化枚举，以总缺口、总峰值配置数、工作量极差和保守资源冗余进行词典序比较；工作量评分沿用该节给定口径，包含架次时长、货箱量、运输能耗和中继服务时长的加权和。枚举规模分别为 $S(13,2)=4095$ 与 $S(13,3)=261{,}625$ 个无标签分区。此处仅转录其结果，不重新枚举。

\subsection{分区方案与缺口比较}
$K=2$ 的最低缺口分区为
\begin{align*}
G_1&=\{S001,S004,S005,S010,S013\},\\
G_2&=\{S002,S003,S006,S007,S008,S009,S011,S012,S014,S015\}.
\end{align*}
合计峰值配置为 A/B/C 型运输机 $1/2/3$ 架、A/B/C 型电池 $1/4/4$ 组、中继无人机 4 架、中继能源组件 4 组。库存分别为 $4/2/2$ 架、$6/4/4$ 组、2 架和 6 组，因此缺口为 C 型运输机 1 架与中继无人机 2 架，共 3 个单位。

$K=3$ 的最低缺口分区为
\begin{align*}
G_1&=\{S001,S004,S005\},\\
G_2&=\{S002,S003,S006,S007,S009,S012,S014,S015\},\\
G_3&=\{S008,S010,S011,S013\}.
\end{align*}
其峰值配置合计为 A/B/C 型运输机 $1/2/3$ 架、A/B/C 型电池 $1/4/5$ 组、中继无人机 4 架、中继能源组件 4 组；相对库存缺少 C 型运输机 1 架、C 型电池 1 组及中继无人机 2 架，共 4 个单位。

\begin{table}[htbp]
\centering\small
\caption{两种分区的峰值配置与库存缺口}\label{tab:q4gap}
\begin{tabular}{c c c c c c c}\toprule
方案 & 运输机 A/B/C & 电池 A/B/C & 中继机 & 中继组件 & 缺口组成 & 缺口合计\\\midrule
$K=2$ & 1/2/3 & 1/4/4 & 4 & 4 & C机1、中继机2 & 3\\
$K=3$ & 1/2/3 & 1/4/5 & 4 & 4 & C机1、C电池1、中继机2 & 4\\\bottomrule
\end{tabular}
\end{table}
图~\ref{fig:q4gapplot} 比较两种方案的缺口结构，图~\ref{fig:q4peakplot} 给出各类资源的峰值配置。$K=2$ 的峰值下界中继机需求为 4，而该更新版小节同时指出固定排程实际使用过的中继实体 ID 为 2；两者分别代表并发下界与排程 ID 统计，不应混为一谈。$K=3$ 将库存拆分到更多独立组后，C 型电池峰值需求增加，因此其总缺口较 $K=2$ 多 1 个单位。
\begin{figure}[htbp]\centering\includegraphics[width=0.82\linewidth]{fig_q4_gap_comparison.png}\caption{两种分区资源缺口}\label{fig:q4gapplot}\end{figure}
\begin{figure}[htbp]\centering\includegraphics[width=0.92\linewidth]{fig_q4_peak_resources.png}\caption{两种分区峰值配置}\label{fig:q4peakplot}\end{figure}
\begin{figure}[htbp]\centering\includegraphics[width=0.62\linewidth]{fig_q4_k2_groups.png}\caption{$K=2$ 分区规模}\label{fig:q4k2}\end{figure}
\begin{figure}[htbp]\centering\includegraphics[width=0.72\linewidth]{fig_q4_k3_groups.png}\caption{$K=3$ 分区规模}\label{fig:q4k3}\end{figure}

本节分区只对冻结的运输和通信任务进行配置核算，不表示问题二、问题三与分区的联合全局最优。新版 Q4 TEX 声明其验证程序复核了服务区恰好覆盖一次和同架次服务区不可拆分；本文沿用该声明，不将之描述为本轮重新运行所得。

\section{敏感性与情景讨论}
\subsection{返航储备参数}
问题一的返航储备敏感性已在图~\ref{fig:q1sens} 中给出。随着预留电量增加，可供载荷使用的能量减少，平均安全载荷呈非增趋势。该趋势与约束式~\eqref{eq:q1payload} 一致；但本文只报告参考结果所含 10\%--30\% 的离散情景，不对区间外进行外推。

\subsection{分区数变化}
问题四的 $K=2$ 与 $K=3$ 对比表明，分区数增加不必然降低总资源需求。独立组不能共享资源，且任务区划改变并发重叠结构；在给定固定排程与库存下，$K=3$ 增加一组 C 型电池缺口，使总缺口从 3 增至 4。此结论针对用户提供的新版 Q4 TEX 所述分区目标与峰值口径，不代表任意时序重排后的总体最优结论。

\subsection{模型边界情景}
若将天气、临时禁飞区或通信衰落引入模型，应相应扩展可行架次集合和链路可用状态；若允许各区之间临时借调资源，则需增加转场时间和跨组调度变量。以上属于后续扩展，本稿不赋予数值预测。

\section{模型评价与推广}
\subsection{模型特点}
本文的分层结构使决策边界可解释：问题一处理单区能力和离散组批，问题二在批次上叠加时限及设备/电池周转，问题三继承运输时间轴配置通信模式，问题四将冻结任务汇总到独立分区核算资源。峰值区间计数与库存缺口有明确的量纲和管理含义，且跨区任务通过连通组件约束保持完整。

\subsection{局限性}
本稿以已有参考解作为报告数据，没有重新运行求解器或独立复核参考解，因此所有数值应按初稿口径使用。问题三的通信模式及中继任务直接沿用参考结果，未在本轮复查各链路输入与实体资源占用。问题四固定运输/通信任务时序，不能解释联合重排后是否会有更小缺口；其峰值并发下界也不等同于考虑位置转场、组间协同后的实体设备规划。绘图所用汇总指标不包含随机天气和设备故障的不确定性。

\subsection{推广方向}
该框架可用于道路受阻地区的应急物资投送、灾时临时通信覆盖和多责任区资源配置。实际部署时，建议以实时气象、道路和通信监测更新输入，并采用滚动时域方式重排剩余任务；同时将中继机位置转场、频谱容量、充电功率与鲁棒备用量纳入统一资源模型。

\subsection{初稿数据来源声明}
问题一至问题三的图表与结果取自工作区 \texttt{results/reference/} 下的归档参考 JSON；问题四数值取自用户提供的更新版 \texttt{sections/8\_problem4.tex}。本文未重新求解 Q3/Q4，亦未审计或独立验证题解。文献检索因本机 SSL 证书校验失败未能完成，故本稿不列未经核实的外部文献。AI 辅助：使用 MathModel 内置 AI 助手（本会话模型标识 GPT-6 Astra）协助整理文本与绘图程序；MathModel/模型发布日期未由当前环境提供，提交前需按实际会话信息补充。

\begin{appendices}
\section{复现材料说明}
本稿结果来自已归档的参考 JSON 和用户提供的新版 Q4 TEX，而非本次执行求解器的输出。图表可由下列命令依据归档 Q1--Q3 JSON 及脚本内明示的新版 Q4 TEX 指标重新生成：
\begin{Python}{论文图表生成命令}
python src/paper_figures.py
\end{Python}
图表脚本将图片输出至 \texttt{figures/}，并使用工作区内的 \texttt{simsun.ttf} 字体。工作区另有旧版参考结果转换脚本，其 Q4 数据仍指向旧归档；本稿未用该脚本生成新版 Q4 结论或表格。此前独立求解脚本也未与本稿引用的归档结果混用。

AI 辅助声明：本文的文字整理与绘图脚本由作者在 MathModel 内置 AI 助手协助下完成。当前会话模型标识为 GPT-6 Astra；具体提供机构配置和发布日期请以提交时会话配置为准，当前环境未提供发布日期。

\end{appendices}
\end{document}
